\label{sec:conclusion}

Algorithm choice has a universally significant effect on training fitness
across all strategy and initialisation conditions, yet this advantage does
not transfer to test performance. The binding constraint is strategy
expressiveness in the face of market non-stationarity: once the target regime
shifts, optimised parameters no longer generalise regardless of how
effectively they were found.

The near-universal training-to-test failure is best attributed to a
misspecified objective rather than defective optimisers. The decisive
evidence is that the train-to-test gap scales with optimisation
\emph{success}, not algorithm identity: GPS, with the lowest training
fitness, has a train/test ratio near $1.0$; SOS, with the highest, shows the
largest gap. Within \texttt{ema\_shared}, training medians span
$\$7{,}900$--$\$128{,}000$ while every algorithm's test median sits at
$\$1{,}000$ (Tables~\ref{tab:perf_train},~\ref{tab:perf_test}). Two
non-exclusive mechanisms underlie this: overfitting to the 2014--2019 price
path, and MA-crossover signals being fundamentally weak predictors across
regimes. The $2.07\%$ baseline-beating rate is consistent with right-tail
sampling rather than a learnable edge --- consistent with, though not a
formal test of, the efficient market hypothesis for short-horizon technical
signals.

\textbf{AOS vs PSO.} AOS performs comparably to PSO on training fitness
despite lacking personal-best memory, suggesting that shell-based population
distribution is sufficient on the lower-dimensional strategies. However, on
\texttt{ema\_independent} AOS collapses on test (median $\$1{,}105$) while
PSO and SOS retain high medians ($\$4{,}435$ and $\$4{,}566$). Since the
primary structural difference is per-agent memory, this suggests personal-best
history aids generalisation on the most expressive landscape --- though the
present design does not isolate memory as a variable.

\textbf{GPS as baseline.} GPS is consistently and significantly worse than
every population-based method on training fitness, confirming that
population diversity is a structural requirement on this multimodal landscape,
not an optional improvement. Its small train-to-test gap reflects weak
optimisation on both splits rather than robustness.

\textbf{SOS compute cost.} SOS's three-evaluation-per-iteration overhead
(wall-clock $\approx 4\times$ PSO) is not justified by its performance: Dunn's
test finds no significant advantage over PSO under matched-iteration budgets.

\textbf{Parameterisation structure.} Independent $\alpha$ in
\texttt{ema\_independent} provides greater expressiveness than shared
$\alpha$, producing higher $H$-statistics and cleaner algorithm separation.
Critically, \texttt{ema\_shared} ($d=3$) transfers worse than the $d=14$
\texttt{weighted} space ($\rho = 0.05$ vs $0.34$): parameterisation
\emph{nature} matters more than size.

\textbf{AOS-specific weakness.} The LE perturbation degrades across
iterations as the LE designation tracks population rank rather than an
absolute landscape property, reducing its escape effectiveness as the nucleus
shifts. This is a genuine weakness in the published formulation.

\textbf{No Free Lunch.} Strategy choice accounts for far more outcome
variance than algorithm choice, consistent with the No Free Lunch
theorem~\cite{wolpert1997}: no algorithm dominates uniformly across
landscapes.

\paragraph*{Outlook}
Future work should adopt risk-adjusted fitness functions (e.g.\ Sharpe ratio
\cite{investopedia_sharpe_ratio}) and walk-forward or cross-validated
evaluation to impose generalisation pressure. Separating the two failure
mechanisms --- objective misspecification vs.\ market non-stationarity ---
requires running the same optimisers against a risk-adjusted objective under
the same split: if the gap narrows, the fitness function dominates; if it
persists, regime shift dominates.